\section{Qiushi Engine and Autonomous Mathematical Research}\label{sec:research}

\subsection{Research Architecture and Computational Environment}

Qiushi Engine is a hierarchical multi-agent research system designed for
sustained autonomous scientific research. Its first
report, on autonomous discovery with a real optical platform, describes
a core research layer for planning, method development, experimentation,
and critical review, supported by knowledge retrieval, memory, auxiliary
exploration, and evidence checking~\cite{yang2026optical}. Curated knowledge
and memory allow earlier findings to inform later decisions without
repeatedly rebuilding the entire context.

Here these capabilities were applied to exact mathematics. Qiushi Engine
carried out the mathematical research end to end: it formulated candidate
approaches, wrote and executed research programs, compared representations,
derived and tested claims, and developed and documented the final proof.
The research initially addressed the broader rank-22 construction and
rank-23 optimality questions. Its subsequent
focus on binary-field bounds concentrated the research on an exact,
finitely representable problem while retaining both constructive and
lower-bound routes. The original scientific report was completed on 9 September.

The computational environment joined Python and C++ programs with
numerical linear algebra, symbolic calculation, combinatorial search,
and proof checking. The released research programs include NumPy and
SymPy calculations and constraint-programming formulations. The final
computational certification uses PySAT encodings, CaDiCaL proof generation, and
\texttt{drat-trim} verification, with a separate C++ reconstruction of
the quotient table. These tools serve different purposes: a numerical
experiment can suggest structure; an exact identity or checked finite
proof establishes it. The maintained mathematical replay runs on Linux
without Qiushi Engine or language-model access. Its dependencies and
commands are collected in \qpath{reproducibility/README.md}. The accompanying
Lean development checks the finite premises through integer branch
certificates and proves the structural deduction within the same formal
framework; its reproduction guide is \qpath{formalization/README.md}.

\subsection{Meta-Trace: Research Memory and Evidence}

Meta-Trace~\cite{yang2026optical} is the structured record of the research process within
Qiushi Engine's memory system. It connects three complementary views:
\emph{history}, preserving summaries of questions, approaches, and actions;
\emph{evidence}, linking claims to arguments, computations, and artifacts;
and \emph{trajectory}, retaining changes in understanding, decisions, and
unresolved questions. Layered summaries and retrieval connect a compact
view of prior research to the detailed records and materials needed to
recover its context. Later reasoning can therefore revisit both a finding
and its conditions, reuse a method, or recognize a previously corrected
error. Meta-Trace provides this research-history foundation for memory;
the underlying evidence, its summaries, and the conclusions drawn from
them remain distinguishable.

A final proof orders statements by logical dependence. Research usually
discovers those dependencies in a different order. A promising candidate
may fail; a weaker formulation may reveal what a stronger one must retain;
an apparent computational obstacle may become a question about algebraic
structure. Understanding this development requires preserving the reasons
for a decision alongside its outcome.

The research trajectory in Appendix~\ref{app:research-record} is a
problem-driven synthesis of the original Meta-Trace records and their
associated notes, plans, programs, and results. Recurring accounts are
condensed; related lines of research are organized by mathematical question;
consequential corrections and unresolved branches remain explicit.
The research guide links these themes to specific supporting artifacts.
This preserves the scientific connections across the research without
reproducing private operational details.

For Qiushi Engine, this record supports continuity and the reuse of
accumulated understanding. For readers, this account makes the
development of the method inspectable: which observation changed a question,
which counterexample invalidated an inference, and which representation
made a new proof possible. It also preserves useful intermediate methods
and open branches as starting points for further research. The final
theorem and the path to it therefore become two connected scientific
objects, each with its own supporting evidence.

Documenting this research trajectory extends the contribution beyond the
final theorem. It makes the development of mathematical understanding
available for study: how evidence changed a conjecture, why a promising
representation was abandoned, and how an intermediate result enabled a
different proof strategy. Researchers can inspect these transitions,
reuse the resulting methods, and compare alternative approaches against
the same evidence. In this way, the open record provides a concrete,
substantive case study of Qiushi Engine's long-horizon autonomous research
capability, expressed through the mathematical decisions and results it
produced rather than through the final answer alone.

Three changes organize the account. Numerical continuation made exactness
and coefficient growth inseparable. Quotient searches exposed the gap
between admissible first-factor supports and complete decompositions.
Finally, equality in a rank bound recovered relations among all three
factors. These are the main points at which accumulated understanding
changed the research strategy. The following subsections explain them;
the appendix develops the intermediate routes and their evidence.

\subsection{Numerical Deformation and Exactness}

Laderman's 23-product construction reorganizes the 27 entrywise products
into shared bilinear work~\cite{laderman1976noncommutative}. Reducing 23 to
22 asks for another exact identity. A 22-term decomposition has
$22(9+9+9)=594$ coefficients and must satisfy 729 cubic coefficient
identities. The small matrix dimensions therefore conceal a highly
coupled algebraic problem.

Qiushi Engine explored deformations of known rank-23 schemes,
pairing-defect formulations, and cyclic and transpose symmetries.
Some numerical residuals decreased while the factor amplitudes grew:
large terms cancelled without approaching a finite-coefficient solution.
This observation changed the question from how far a residual could be
reduced to what geometry would support an exact shorter decomposition.
Incidence formulations and local compatibility tests made that distinction
explicit. They supplied reusable representations and checks while the
rank-22 construction question remained open.

\subsection{Quotient Cores and Factor Compatibility}

Over $\FF$, exact arithmetic removes numerical zero-testing from the
existence problem. A codimension-one core obtained by quotienting the
$E_{11}$ direction offered a concrete target: a 19-term decomposition of
that core, embedded in the coordinate complement and supplemented by the
three deleted products, would yield a full 22-term decomposition.

Occupation bounds constrained the possible first factors of such a
candidate. Yet an admissible first-factor multiset did not supply
compatible second and third factors. Qiushi Engine examined fixed-factor
completion, contraction ranks, lift compatibility, coding interpretations,
and moment relaxations to understand this missing information. Each
representation offered a different way to test compatibility; several
necessary-condition models remained feasible or computationally unresolved.

The lasting insight was the distinction between locating factors and
coupling them. It motivated a return to the full tensor: could support
geometry force a configuration in which additional algebraic identities
become unavoidable, without first deciding the entire core problem?

\begin{figure}[tbp]
  \centering
  \makebox[\linewidth][c]{\begin{tikzpicture}[x=1cm,y=1cm,
  route/.style={draw=ReportRule,fill=white,align=center,text width=6.2cm,
    inner sep=8pt,font=\sffamily\small,minimum height=1.45cm},
  flow/.style={-{Stealth[length=2.2mm]},thick,draw=ReportPrimary}]
  \node[route] (broad) at (0,0)
    {\textbf{Exact algorithm search}\\Known rank-23 schemes;\\deformation and symmetry};
  \node[route] (core) at (7.15,0)
    {\textbf{Finite-field restrictions}\\Core supports; occupation;\\factor completion};
  \node[route] (lesson) at (0,-2.15)
    {\textbf{Mathematical lessons}\\Finite exactness;\\support versus factor compatibility};
  \node[route] (full) at (7.15,-2.15)
    {\textbf{Return to the full tensor}\\High-rank pairs; affine cosets;\\rank-sum equality};
  \node[route,text width=13.35cm,draw=ProofGreen] (result) at (3.575,-4.3)
    {\textbf{Structural lower-bound proof over $\FF$}\\Eight certified quotient raises
    $\longrightarrow$ affine profile $(16,1,3)$\\
    Saturation $\longrightarrow$ cross-factor identities $\longrightarrow$ contradiction};
  \draw[flow] (broad) -- (lesson);
  \draw[flow] (core) -- (full);
  \draw[flow] (lesson.east) -- (full.west);
  \draw[flow] (full.south) -- (full.south |- result.north);
\end{tikzpicture}
\unskip}
  \caption{Changes of mathematical representation during the research.
  Construction and quotient searches clarified exactness and factor
  compatibility. Returning to the full tensor connected these questions
  to a rank equality and the structural lower-bound proof.}
  \label{fig:research-route}
\end{figure}

\subsection{Split-Flattening Saturation and Algebraic Structure}

For a hypothetical 20-term decomposition, the split flattening requires
$\sum_t\rk(A_t)\ge27$. Certified occupation bounds and affine geometry
force sixteen rank-one, one rank-two, and three rank-three first factors.
Their rank sum is $16+2+9=27$: the inequality is an equality.

Recognizing this profile statistic as a rank-additivity condition changed
the final proof obligation:
\[
  P=\sum_t M_t,\qquad \sum_t\rk(M_t)=\rk(P)
  \quad\Longrightarrow\quad M_tP^{-1}M_s=\delta_{ts}M_t.
\]
Equality couples the second and third factors through $P^{-1}$. For the
matrix-multiplication tensor, evaluating these products explicitly makes
the required invertible factors incompatible. Thus first-factor counting
leads to a contradiction involving all three factors. The final argument
uses finite geometry to force equality and equality to impose algebraic
structure, replacing the remaining support enumeration with an identity.

Qiushi Engine also reduced the finite assumptions. Only eight
two-dimensional orbits, in which all three nonzero matrices have rank
at least two, are needed. An additional orbit considered earlier already
has a rank-one third direction. Removing it shortened the proof by
identifying the precise statement the geometry required.

\subsection{Semantic Validation and Independent Checks}

The finite premises were reconstructed through separate encodings:
bounded Boolean copies of integer multiplicities and explicit monotone
unary counts. One quotient allows a direction of multiplicity two, so
the distinction between integer counts and distinct-support bits is
essential. An earlier auxiliary-variable collision had exposed a false
contradiction. The corrected approach therefore supplied an explicit
mathematical-to-Boolean assignment map as well as a checked refutation.

The cross-factor identity received a similarly direct test. The actual
$27\times27$ permutation matrix was constructed and the products of
elementary matrices were computed independently of the proposed index
formula. The recorded checks cover all $729^2=531{,}441$ pairs of
elementary simple tensors, complementing the multilinear derivation
below. A convention error encountered along the way was corrected at
the level of the tensor identity.

These checks changed the mathematical explanation: they made multiplicity,
coordinates, and the scope of each implication explicit. The resulting
proof has a finite computational foundation followed by a fully written
structural deduction. The completed Lean formalization connects both
parts to the actual matrix-multiplication semantics: the finite premises
are proved, and the algebraic argument yields a closed theorem for every
input matrix. Section~\ref{sec:lean} records its statement coverage and
kernel checks.

\subsection{Research Development and Reusable Knowledge}

The proof and the Meta-Trace answer different scientific questions.
The proof establishes why a 20-term decomposition is impossible.
The Meta-Trace explains how Qiushi Engine arrived at the representation
that made this contradiction accessible. Early work contributed exact
controls, quotient constructions, and an understanding of the information
lost by relaxations; later work used those distinctions to formulate a
more effective question about the full tensor.

Together, the account and its research materials expose the development
of a method: which ideas were proposed, what calculations tested them,
why some interpretations were revised, and how retained understanding
informed subsequent research. This makes the autonomous research
process itself available for scientific study and reuse, alongside its
mathematical result.

\begin{figure}[tbp]
  \centering
  \makebox[\linewidth][c]{\begin{tikzpicture}[x=1cm,y=1cm,
  box/.style={draw=ReportRule,fill=white,align=center,text width=6.2cm,
    inner sep=8pt,font=\sffamily\small,minimum height=1.45cm},
  flow/.style={-{Stealth[length=2.2mm]},thick,draw=ReportPrimary}]
  \node[box] (finite) at (0,0)
    {\textbf{Certified finite premises}\\Line bounds; eight 2-plane raises;\\affine-plane bounds};
  \node[box] (split) at (7.15,0)
    {\textbf{Full-tensor split flattening}\\$\sum_t\rk(A_t)\ge\rk(P)=27$\\At least four high-rank first factors};
  \draw[dashed,ReportRule] (-3.38,-1.08) -- (10.53,-1.08);
  \node[box,text width=13.35cm] (geometry) at (3.575,-2.15)
    {\textbf{High-rank factors share an affine row or column coset}\\
     Pairwise differences have rank one;\\
     the normalized lower block has rank two};
  \node[box,text width=13.35cm] (profile) at (3.575,-4.3)
    {\textbf{Affine-plane caps force the profile $(16,1,3)$}\\
     Exactly four high-rank first factors, three invertible\\
     $\sum_t\rk(A_t)=16+2+9=27$: equality in the rank bound};
  \node[box,text width=13.35cm] (sat) at (3.575,-6.45)
    {\textbf{Equality forces cross-factor identities}\\
     $M_tP^{-1}M_s=\delta_{ts}M_t$\\
     $M_tP^{-1}M_s=M(A_tB_sC_t^TA_s,B_t,C_s)$};
  \node[box,text width=13.35cm,draw=ProofRed] (end) at (3.575,-8.6)
    {\textbf{Contradiction: at most one invertible first factor, but three are required}\\
     Diagonal identities force $B_t,C_t$ to be invertible when $A_t$ is invertible\\
     Off-diagonal identities would make an invertible product zero};
  \draw[flow] (finite.south) -- (finite.south |- geometry.north);
  \draw[flow] (split.south) -- (split.south |- geometry.north);
  \draw[flow] (geometry) -- (profile);
  \draw[flow] (profile) -- (sat);
  \draw[flow] (sat) -- (end);
\end{tikzpicture}
\unskip}
  \caption{Two complementary constraints meet at equality.
  The finite-input branch restricts the first-factor geometry; the
  flattening branch supplies a rank budget. At the forced profile the
  budget is exactly exhausted, enabling a symbolic contradiction involving
  all three factors. Above the dashed line are the finite premises and
  the general rank bound; below it are their structural consequences
  for a hypothetical 20-term decomposition over $\FF$.}
  \label{fig:proof-structure}
\end{figure}
\FloatBarrier
